\chapter{Immunisations and Vaccinations}
\index{Immunisations and Vaccinations}

\section*{Definition and Overview}
Immunisation is the process by which a person becomes protected against an infectious disease through vaccination. A vaccine contains a safe form of a pathogen, or a piece of it, that trains the immune system to recognise and fight the real infection without causing the disease. some countries have one of the most successful immunisation programs in the world, protecting people of all ages against a wide range of serious diseases, from measles and whooping cough in childhood to influenza, shingles, and pneumococcal disease in older adults. Vaccines are among the most rigorously tested and effective public health measures, saving millions of lives globally each year.

\section*{Epidemiology}
Vaccination has transformed the health of people. Diseases that once killed or disabled thousands of children, such as polio, measles, diphtheria, and tetanus, are now rare or eliminated, and the rates of childhood vaccination in many countries are high, with around 95\% of two-year-olds fully vaccinated. Influenza vaccination is recommended annually for all people aged six months and over, and specific vaccines are funded for older adults, pregnant women, Aboriginal and Torres Strait Islander people, and people with certain medical conditions. When vaccination rates fall, outbreaks of vaccine-preventable diseases return, as has occurred with measles in some communities, underscoring the importance of maintaining high coverage.

\section*{Aetiology and Pathophysiology}
Vaccines work by presenting the immune system with antigens, the molecules of a pathogen that the immune system recognises. Different vaccine technologies achieve this in different ways: live attenuated vaccines use weakened versions of the whole organism; inactivated vaccines use killed organisms; subunit and conjugate vaccines use specific proteins or sugars; and newer mRNA vaccines carry instructions for the body's own cells to produce a harmless piece of the pathogen. In all cases, the immune system responds by producing antibodies and memory cells. On later exposure to the real pathogen, the memory response is rapid and strong, preventing disease or greatly reducing its severity. Herd immunity adds a further layer of protection: when enough of the population is vaccinated, the spread of infection is interrupted, protecting those who cannot be vaccinated.

\section*{Clinical Features}
\begin{itemize}
    \item Childhood schedule: hepatitis B at birth, then a schedule protecting against diphtheria, tetanus, whooping cough, polio, Hib, pneumococcal, rotavirus, measles, mumps, rubella, meningococcal, and varicella diseases
    \item School-based programs: HPV and some meningococcal vaccines for adolescents
    \item Adult vaccines: influenza annually, COVID-19, and boosters for tetanus and whooping cough
    \item Vaccines for older adults: shingles, pneumococcal, influenza, and COVID-19
    \item Vaccines for special groups: hepatitis A and B, travel vaccines, and vaccines for people with medical conditions
    \item Common reactions are mild: sore arm, low-grade fever, and tiredness
    \item Protection may require multiple doses for full immunity, and some vaccines need boosters
\end{itemize}

\section*{Differential Diagnosis}
Not every infection is preventable by vaccination, and individual recommendations vary.
\begin{itemize}
    \item \textbf{Vaccines available but not funded:} some vaccines, such as for Japanese encephalitis or hepatitis A, are recommended but must be purchased
    \item \textbf{Travel-specific vaccines:} yellow fever, typhoid, and cholera depend on destination and itinerary
    \item \textbf{Medically indicated vaccines:} for people with chronic disease, immunosuppression, or occupational risk
    \item \textbf{Catch-up vaccination:} for people who have missed doses at any age
    \item \textbf{Immunisation for immunity:} some people have natural immunity from prior infection and may not need certain doses
\end{itemize}

\section*{When to Seek Medical Care}
Everyone should have a regular review of their immunisation status, at least when they visit a GP for a health check, during pregnancy, and before overseas travel. People starting immunosuppressive treatment should have vaccination completed before treatment begins. Anyone with a medical condition, planned pregnancy, or travel plans should discuss their vaccine needs well in advance.

\textbf{Contact your local emergency services for:}
\begin{itemize}
    \item Difficulty breathing, wheezing, or swelling of the face or throat within minutes of a vaccination (anaphylaxis)
    \item Collapse or a seizure after vaccination
    \item Signs of a severe allergic reaction with a rapidly spreading rash
\end{itemize}

\section*{Diagnosis and Investigations}
\begin{itemize}
    \item \textbf{Immunisation history review:} checking the local Immunisation Register (AIR) record
    \item \textbf{Serology:} blood tests to check immunity, for example for hepatitis B, measles, or rubella in specific situations
    \item \textbf{Risk assessment:} age, medical conditions, pregnancy, occupation, and travel plans determine recommendations
    \item \textbf{Schedule planning:} catch-up schedules for missed doses, planned around other treatments
\end{itemize}

\section*{Management}
\begin{itemize}
    \item \textbf{Childhood immunisation:} follow the National Immunisation Program schedule from birth to adolescence
    \item \textbf{Adult boosters:} tetanus, diphtheria, and whooping cough booster every 10 years in adulthood
    \item \textbf{Annual influenza:} recommended for everyone aged six months and over
    \item \textbf{COVID-19:} primary course and boosters as advised by ATAGI
    \item \textbf{Vaccination in pregnancy:} whooping cough and influenza in every pregnancy
    \item \textbf{Special risk groups:} pneumococcal and shingles vaccines for older adults and those with chronic disease
    \item \textbf{Travel vaccination:} planned 6--8 weeks before departure
    \item \textbf{Catch-up vaccination:} available at any age through GPs, with free catch-up for some vaccines under the National Immunisation Program
\end{itemize}

\section*{Complications}
\begin{itemize}
    \item Minor reactions: soreness, redness, fever, and fatigue, which settle within days
    \item Very rare serious reactions, including anaphylaxis, occurring in about one in a million doses
    \item Vaccine-preventable diseases when vaccination is missed: measles, whooping cough, influenza, and pneumococcal disease can all cause serious illness, hospitalisation, and death
    \item Complications of disease in pregnancy, older age, and chronic illness make vaccination particularly important in these groups
\end{itemize}

\section*{Prognosis and Outcomes}
Fully vaccinated people are protected against the targeted diseases at very high rates, and the population-level effects of immunisation have been profound: smallpox is eradicated, polio is eliminated in many countries, and measles, rubella, and Hib disease have been reduced to very low levels. Immunity from some vaccines wanes over time, which is why boosters and annual influenza vaccination matter. High vaccination coverage protects the whole community, including the most vulnerable, and is one of the best investments in health available.

\section*{Prevention and Self-Care}
\begin{itemize}
    \item Keep your immunisation records up to date, including on the local Immunisation Register
    \item Book childhood vaccinations on time according to the schedule
    \item Ask about vaccines during health checks, pregnancy, and before travel
    \item Have the influenza vaccine every year before winter
    \item Ensure the whole family is vaccinated to protect young babies and vulnerable people
    \item Get reliable information from your GP or government health sources rather than social media
    \item Report any significant reactions to your doctor
\end{itemize}

\section*{Sources and Further Reading}
The following reputable sources provide further information for healthcare students and patients seeking reliable, current advice about immunisation and vaccination.
\begin{itemize}
    \item Australian Government Department of Health and Aged Care. \textit{National Immunisation Program.} https://www.health.gov.au/
    \item National Centre for Immunisation Research and Surveillance. \textit{Vaccine information for the public.} https://ncirs.org.au/
    \item Healthdirect Australia. \textit{Immunisation and vaccines.} https://www.healthdirect.gov.au/
    \item World Health Organization. \textit{Vaccines and immunisation.} https://www.who.int/
    \item NHS. \textit{Vaccinations.} https://www.nhs.uk/vaccinations/
    \item Centers for Disease Control and Prevention. \textit{Vaccines and immunizations.} https://www.cdc.gov/vaccines/
\end{itemize}
